\documentclass[11pt]{article}
\usepackage[T1]{fontenc}
\usepackage{lmodern,microtype}
\usepackage{amsmath,amssymb,amsthm}
\usepackage[letterpaper,margin=1in]{geometry}
\usepackage{booktabs,array,enumitem}
\usepackage{xcolor}
\usepackage{listings}
\usepackage[colorlinks=true,linkcolor=blue!50!black,citecolor=blue!50!black,urlcolor=blue!50!black]{hyperref}
\lstset{basicstyle=\ttfamily\small,columns=fullflexible,frame=single}
\newcommand{\Q}{\mathbb Q}
\newcommand{\A}{\mathbb A}
\newcommand{\C}{\mathbb C}
\newcommand{\F}{\mathbb F}
\newcommand{\Z}{\mathbb Z}
\newcommand{\GL}{\mathrm{GL}}
\newcommand{\p}{\mathfrak p}
\newcommand{\Reg}{\operatorname{Reg}}
\newcommand{\Ptag}{\textbf{(P)}}
\newcommand{\Etag}{\textbf{(E)}}
\newcommand{\Ctag}{\textbf{(C)}}
\theoremstyle{definition}
\newtheorem{fact}{Fact}
\title{The running example for the metaplectic Jacquet--Langlands book:\\
\normalsize $K=\Q(\zeta_3)$, $n=3$, $r=2$, $D=(-1,\,4-\zeta_3)_K$}
\author{David Joyner \and Claude (Anthropic)}
\date{November 8, 2026 --- draft; resynchronized to the book November 9, 2026}
\begin{document}
\maketitle

\begin{abstract}
We fix one concrete case of the standing assumptions of \cite{Book26} and
\cite{JC26}, and evaluate in it the main objects of each chapter. These are the
cocycle parameters, the commutator pairings, the good elements, the
constants $b_D$ and $c_D$, the wild toral decompositions, and the global
normalization ledger. Every statement is either checked here, or quoted
with its location and status label. The aim is to make the general
constructions concrete, and to give each chapter a common place to test its
formulas.
\end{abstract}

\paragraph{Theorem numbers.} The four parts and the companion paper on
the rank-two wild place have been combined into a single book
\cite{Book26}, which is now the current source; references below are
to its chapter, section and theorem numbers (Chapters~4--7 are Parts~I--IV,
and Chapter~8 is the companion paper). Numbers are current as of the
book draft committed alongside this revision and can shift in later
drafts, since the book is renumbered automatically when sections are
added or removed.

\section{The data}
\begin{center}
\begin{tabular}{@{}ll@{}}
\toprule
Base field & $K=\Q(\zeta_3)$, totally imaginary, $\mu_3\subset K$\\
Covering degree, rank & $n=3$, \quad $r=2$, \quad $(n,r)=1$\\
Parameters (Assumption~4.1.1) & $r'=2$, \quad $c=1$, \quad $c_v=r'c=2$\\
Checks & $3\mid r-1+2c=3$, \quad $3\mid r-1+2rc_v=9$\\
Division algebra & $D=(-1,\,4-\zeta_3)_K$, \quad $4-\zeta_3=(1-\zeta_3)(3+\zeta_3)$\\
Ramification & $S=\{\p_3,\p_7\}$, \quad $\p_3=(1-\zeta_3)$, \quad $\p_7=(3+\zeta_3)$\\
Cyclic subfield (\cite{JC26}) & $L=K(i)=\Q(\zeta_{12})\subset D$\\
\bottomrule
\end{tabular}
\end{center}

\begin{fact}[\Ptag\ Ramification]
$D$ is ramified exactly at $\p_3$ and $\p_7$.
\end{fact}
\begin{proof}
\emph{At $\p_3$ and $\p_7$.} The element $-1$ is a nonsquare in the
residue fields $\F_3$ and $\F_7$, since $3\equiv7\equiv3\pmod4$. The
element $4-\zeta_3$ has valuation~$1$ at each of these places. So
$(-1,4-\zeta_3)=-1$ at both.

\emph{At the other finite odd places.} Both entries are units, so the
symbol is~$1$. This includes the other prime above~$7$, namely
$(3+\zeta_3^2)$, since $7$ splits in $K$.

\emph{At~$2$.} The product formula forces the symbol to be~$1$.

\emph{Machine check.} Sage gives:
\begin{lstlisting}
K.<z> = CyclotomicField(3); a, b = K(-1), 4 - z
for p in [2,3,5,7,13]:
    for P, e in K.ideal(p).factor():
        print(p, P, K.hilbert_symbol(a, b, P))
\end{lstlisting}
This returns $-1$ exactly at $(z-1)$ and at $(-3z-2)$, and $+1$ at
$(2)$, $(5)$, $(3z+1)$, $(z+4)$ and $(z-3)$. Here $(-3z-2)=(3+z)$,
since $(-3z-2)/(3+z)=z^2$ is a unit.
\end{proof}

The hypotheses of Assumption~4.1.1 (Chapter~4, continuing those of
\cite[\S1]{JC26}) also hold.
\begin{itemize}
\item $L=K(i)$ is cyclic of degree $r=2$ and lies in $D$.
\item At both places of $S$, $L_v$ is a field, because $-1$ is not a
square in $\F_3$ or $\F_7$.
\item $|S|=2$ is even, as it must be for $r=2$.
\end{itemize}
A fuller check of this example against the hypotheses of \cite{JC26},
\cite{FK86} and \cite{FK88} is in the book's Chapter~1
(Example~1.3.1 and Remark~1.3.2).

\section{The local fields}
\begin{center}
\begin{tabular}{@{}lllll@{}}
\toprule
Place & $K_v$ & Type for $n=3$ & $|3|_v$ & $|K_v^\times/K_v^{\times3}|$\\\midrule
$\p_3\in S$ & $\Q_3(\zeta_3)$ ($e=2$, $q=3$) & wild & $1/9$ & $81$\\
$\p_7\in S$ & $\Q_7$ (contains $\mu_3$) & tame & $1$ & $9$\\
$(2)$ & $\Q_2(\zeta_3)$ ($f=2$, $q=4$) & tame, split for $D$ & $1$ & $9$\\
$(3+\zeta_3^2)$, $(13)=\p\p'$, \ldots & $\Q_\ell$, $\ell\equiv1\ (3)$ & tame, split & $1$ & $9$\\
complex & $\C$ & cover splits & $9$ & $1$\\
\bottomrule
\end{tabular}
\end{center}
The last column is $n^2|n|_v^{-1}$ at the finite places (Chapter~7,
Proposition~7.4.2). Locally at $\p_3$ we have $3+\zeta_3\equiv1\pmod{\p_3}$,
which is a square. So $D_{\p_3}\cong(-1,\varpi)$ with
$\varpi=\zeta_3-1$; this is exactly the model $E_0\oplus E_0j$, $j^2=\varpi$,
of the companion chapter, Chapter~8, Example~8.7.1. At $\p_7$, $D_{\p_7}\cong(-1,7)_{\Q_7}$.

\section{Chapter 4: cocycle, commutators, good elements}
\begin{itemize}
\item \emph{Diagonal torus} (Chapter~7, Proposition~7.2.1). With $c_v=2$,
\[
[\tilde t,\tilde u]=\prod_{i,j}(t_i,u_j)_3^{A_{ij}},\qquad
A=-I+5J\equiv\begin{pmatrix}1&-1\\-1&1\end{pmatrix}\pmod 3,
\]
so $[\tilde t,\tilde u]=(t_1/t_2,u_1/u_2)_3$. A regular
$\mathrm{diag}(t_1,t_2)$ is good if and only if $t_1/t_2$ is a cube.
\item \emph{Quadratic tori} (Chapter~4, Proposition~4.3.1 and Example~4.3.7).
Here $[x,y]=(Nx,Ny)_F^{2}(x,y)_E^{-1}$. The good elements are
$F^\times E^{\times3}$ (Chapter~4, Proposition~4.3.5).
\item \emph{Index of the good elements in a quadratic torus},
$|3|_v^{-1}$ (Chapter~4, Proposition~4.3.5):
\begin{itemize}
\item at $\p_3$ it is $9$: the covered torus is a Heisenberg group, and
its genuine irreducibles have dimension~$3$;
\item at $\p_7$ it is $1$: every element is good, and the covered torus
is abelian.
\end{itemize}
\item \emph{Sign.} The matching sign $(-1)^{r-1}$ of Chapter~4,
Theorem~4.4.5 (Theorem~B), is $-1$.
\item \emph{A regular-but-not-good element} (Chapter~7, Remark~7.3.3). The
element $\mathrm{diag}(1,\zeta_3)$ is regular. It is not good at
$\p_3$, because $\zeta_9\notin\Q_3(\zeta_3)$. Its cube is the identity.
\end{itemize}

\section{Chapter 6: the transform $L_D$}
\begin{center}
\begin{tabular}{@{}lccc@{}}
\toprule
& $b_D=|3|^{-1/2}$ & $c_D=|3|^{-3/2}$ & $R_E=\Delta(t)/\Delta(t^3)$\\\midrule
$\p_7$ (tame) & $1$ & $1$ & $\equiv1$ (Chapter~6, Lemma~6.7.4)\\
$\p_3$ (wild) & $3$ & $27$ & $\in\{1,3,9\}$ (below)\\
\bottomrule
\end{tabular}
\end{center}

\begin{itemize}
\item \emph{At $\p_7$.} $L_D$ is a dimension-preserving bijection of
irreducible characters (Chapter~6, Theorem~6.7.6) \Ptag.
\item \emph{At $\p_3$, divisibility.} The cover does not split, and
every genuine irreducible has dimension divisible by~$3$ (Chapter~6,
Proposition~6.8.1) \Ptag.
\item \emph{At $\p_3$, the prediction.} If $L_D$ maps irreducibles to
irreducibles, then $\dim\tau=27\dim\rho$. In particular $L_D\mathbf1$
would be an irreducible of dimension~$27$.
\end{itemize}

\section{Chapters 7 and 8: the wild place $\p_3$}
\subsection*{The three quadratic tori}
$F=\Q_3(\zeta_3)$ has three quadratic extensions. All of them embed in
$D_{\p_3}$. On $E^1\cong E^\times/F^\times$, Chapter~7's
toral-positivity analysis gives the following \Ptag.

\begin{center}
\begin{tabular}{@{}llll@{}}
\toprule
$E/F$ & values of $R_E$ & $R_E$ as a permutation character & coefficients\\\midrule
unramified $F(i)$ & $9,\ 1$ & $1+2\Reg(C_4)$ & $(3,2,2,2)$\\
ramified $F(\sqrt{u\varpi})$, either type & $9,\ 3,\ 1$ & $1+\Reg(C_2)+\Reg(C_6)$ & $(3,2,1,1,1,1)$\\
\bottomrule
\end{tabular}
\end{center}

Both expansions have value $9=|3|_F^{-1}$ at the identity. So the
restriction of $L_D\mathbf1$ to every covered quadratic torus is an
actual genuine character of dimension $3\cdot9=27$, consistent with the
prediction above. In the cyclic algebra model $j^2=\varpi$:
\begin{center}
\begin{tabular}{@{}cccc@{}}
\toprule
$t$ & $j$ & $1+j$ & $1+\varpi j$\\\midrule
$R(t)$ & $1$ & $3$ & $9$\\
$(L_D\mathbf1)(t^*)$ & $3$ & $9$ & $27$\\
\bottomrule
\end{tabular}
\end{center}

\subsection*{What the book now proves}
This question is no longer open for the running example. Chapter~8
(the companion chapter) proves, as its main theorem (Target~A) \Ptag:
$L_D\mathbf1=\Theta_\tau$ for a unique
genuine irreducible representation $\tau$ of $\widetilde{D_{\p_3}^\times}$,
of dimension $27$, exhibited explicitly as induced from a canonical
genuine character of the third congruence subgroup. The proof runs
through hypothesis (H1) (Chapter~8, Proposition~8.11.4), the theorem
that (H1) implies Target~A (Chapter~8, Theorem~8.10.6), and the
uniqueness/irreducibility of the restriction to the pro-$3$ part
(Chapter~8, Theorem~8.9.5). The same chapter extends this to every
representation of $D^\times$ trivial on the third congruence subgroup
and its norm twists. What remains open -- rows of higher level in
this example, and the wild place for other primes and base fields --
is discussed in the book's Chapter~8, \S8.15 (Future directions).
This is no longer the smallest open instance of the wild integrality
problem; it is the first solved one.

\section{Global: good classes and the normalization ledger}
\begin{itemize}
\item \emph{Good classes} (Chapter~7, Theorem~7.5.2) \Ptag. A regular
$\gamma\in D^\times(K)$ with $E=K(\gamma)$ has a rational lift that is
good everywhere if and only if $\gamma\in K^\times E^{\times3}$. The map
$[s]\mapsto[s^3]$ is a bijection onto these classes.
\item \emph{The ledger} (Chapter~7, Example~7.8.8). Take $\lambda_v=1$ at the
finite split places. Then
\[
\Lambda=\lambda_\infty\cdot b_{D,\p_3}\cdot b_{D,\p_7}=3\lambda_\infty,
\]
and $\Lambda=1$ forces $\lambda_\infty=1/3=|3|_\infty^{-1/2}$.
\item \emph{The identity-term heuristic.} It gives the complex factor
$\lambda_\infty|3|_\infty^{-1}=1/27$, which balances $c_{D,\p_3}=27$ by
the product formula.
\item \emph{Unramified places.} \cite{FK86}'s unit matching needs $p\nmid n$ and
$(n,N_2)=(3,1)=1$. Both hold at every finite place outside $S$, since $\p_3\in S$
is the only place above~$3$. (As the book's Chapter~1, Remark~1.3.2
notes, $(n,N_2)=1$ holds automatically for every odd $n$ in the
quaternion case $r=2$, since there are no composite integers at
most~$2$.)
\end{itemize}

\section{Relation to \cite{JC26}}
In \cite{JC26} the lifting runs from genuine automorphic representations of
an $n$-fold cover of $L_{\A}^\times$, here with $L=\Q(\zeta_{12})$, to
$\varepsilon$-invariant genuine representations of the cover of
$D^\times_{\A}$. Chapter~4, \S4.1, of the book says this. A section of
\cite{JC26} evaluated in this example would complete the picture.

\section*{Suggested use}
Each chapter of the book has a short subsection, ``A running example''
(or, in Chapter~1, the canonical statement, Example~1.3.1), pointing
to the relevant section of this note; \cite{Book26} cites this note
in the normal way wherever the example is worked out further. Each
new formula in a draft can be tested here first. Chapter~8,
\S8.15 is the natural place to look for the next target in the wild
integrality problem, now that this example's own wild place is
settled.

\begin{thebibliography}{99}
\bibitem{Book26} D.~Joyner, Claude and ChatGPT, \emph{On a metaplectic
Jacquet--Langlands transfer for division algebras of prime degree},
book draft, \texttt{book/mjl\_book\_draft.tex}, this repository, 2026.
\bibitem{JC26} D.~Joyner and Claude, On the metaplectic analog of
Kazhdan's ``endoscopic'' lifting (original 1988, corrected 2026),
\texttt{metaplectic\_kazhdan\_revised12.tex}.
\bibitem{FK86} Y.~Z.~Flicker and D.~A.~Kazhdan, Metaplectic
correspondence, \emph{Publ.\ Math.\ IH\'ES} 64 (1986), 53--110.
\bibitem{FK88} Y.~Z.~Flicker and D.~A.~Kazhdan, A simple trace
formula, \emph{J.\ Analyse Math.} 50 (1988), 189--200.
\end{thebibliography}

\end{document}
